\section{The remaining real comparison obstruction before \texorpdfstring{\((\mathcal R_B)\)}{the real-zero hypothesis}}
\label{sec:s5-abscissa-reduction}

The Fourier-rigidity theorem can be used before assuming
\((\mathcal R_B)\) to exclude an entire convergence branch.
It does not, by itself, put the divisor of a remaining comparison
on the critical line.  We prove a sharper unconditional reduction
of that remaining branch and identify precisely what its boundary
positivity supplies.

Retain the specified Calegari family and the notation
\(A=L_f(\,\cdot\,,\varpi)\), \(C=L(\,\cdot\,,\eta)\),
\(B=C L_W\), \(L_V,L_{V\eta},L_X\) from MA-S5.
The gamma factors of \(A\) and \(L_W\) are exactly equal,
as proved in MA-FQ-001.  Write
\[
 R(s)=\frac{A(s)}{L_W(s)}
       =U_{\mathcal B}(s)D_\Omega(s),\qquad
 D_\Omega(s)=\prod_{p\in\Omega}
                  \left(\frac{1+p^{-s}}{1-p^{-s}}\right)^2 .
 \tag{MA-RB.17}
\]
The finite factor \(U_{\mathcal B}\) is holomorphic and nonzero
on \(\operatorname{Re}s>0\).  Indeed every finite local parameter
in this comparison has finite image and zero monodromy by
the checked finite-base-change argument in MA-S5-003;
its local roots have modulus at most one.
The omitted archimedean ratio is one.
Conductor powers in the fully completed quotient are
nowhere-zero exponentials, so they change none of its orders.

\paragraph{MA-RB-002 (a shared real zero and multiplicity at least nine).}
If the specified all-place comparison does not hold, then
\(\Omega\) is infinite and the nonnegative Dirichlet series of
\(D_\Omega\) has a finite abscissa \(\gamma\) with
\[
                 \tfrac12<\gamma<1.
 \tag{MA-RB.18}
\]
It has a real pole there.  Put
\[
 \begin{gathered}
 k=-\operatorname{ord}_\gamma R
     =-\operatorname{ord}_\gamma D_\Omega\in\mathbb Z_{\ge1},
 \qquad n=\operatorname{ord}_\gamma A,\quad
 b=\operatorname{ord}_\gamma C,\\
 v_0=\operatorname{ord}_\gamma L_V,\quad
 v_1=\operatorname{ord}_\gamma L_{V\eta},\quad
 e=\operatorname{ord}_\gamma L_X .
 \end{gathered}
 \tag{MA-RB.19}
\]
Then
\[
 \begin{aligned}
 \operatorname{ord}_\gamma L_W&=n+k,\\
 \operatorname{ord}_\gamma B&=b+n+k,\\
 v_1+e&\ge k,\\
 \operatorname{ord}_\gamma\zeta_L&\ge b+5n+9k\ge9k .
 \end{aligned}
 \tag{MA-RB.20}
\]
Thus \(B\) and \(L_{V\eta}L_X\) have a common real zero at
\(\gamma\).  With \(\mathcal H_L=1+\frac12\log\Delta_L\)
as in MA-RB-001, one has the sharper bound
\[
       1-\gamma\ \ge\
                \frac{(3\sqrt{k}-1)^2}{\mathcal H_L}
                       \ \ge\ \frac4{\mathcal H_L}.
 \tag{MA-RB.21}
\]
These conclusions use neither \(\mathsf A\), \(M\), nor
\((\mathcal R_B)\).

In particular, absence of a common real zero of
\(B\) and \(L_{V\eta}L_X\) in
\[
                 \left(\tfrac12,\,1-\frac4{\mathcal H_L}\right]
 \tag{MA-RB.22}
\]
suffices for the exact comparison.
An empty interval satisfies that sufficient condition.
This is a comparison criterion; it is not a proof of
\((\mathcal R_B)\) or a zero-free assertion about every
zero of \(B\).

\begin{proof}
\emph{The convergence branch excluded by Fourier rigidity.}
The classical finite-\(\Omega\) comparison, Calegari's
Lemma~4.10, gives exact equality at all places.
Consequently failure implies \(\Omega\) is infinite.
Each Euler factor in (MA-RB.17) has nonnegative Taylor
coefficients, and its coefficient of \(p^{-s}\) is four.
Their product therefore gives a nonnegative Dirichlet
series in the absolute-convergence half-plane.

The previously checked proof of Calegari's Theorem~4.7
gives \(\sum_{p\in\Omega}p^{-1}<\infty\).
The Euler product consequently has a finite positive
limit at \(s=1\) from the right.
Its meromorphic continuation is (MA-RB.17).
If its abscissa were one, Landau's theorem would force
a singularity at one; meromorphy and that finite positive
limit make the point regular instead.
Thus its abscissa is less than one.
It is finite below: any one prime in \(\Omega\) already
has a nonnegative local series diverging on the real
half-line \(s\le0\).
Hence \(\gamma\ge0\).

Suppose \(\gamma\le1/2\).
For each real \(\sigma>1/2\), convergence of the positive
Dirichlet series and its prime coefficient four give
\(\sum_{p\in\Omega}p^{-\sigma}<\infty\).
The Euler logarithm then converges absolutely and locally
uniformly on \(\operatorname{Re}s>1/2\).
Together with the finite local unit this makes \(R\)
holomorphic and nonzero throughout that half-plane.

The genuine comparison \(L\)-datum, its integral signed
orders, its exact zero gamma kernel, its prime-power
support and eventual nonnegative good-prime coefficients
were constructed in MA-FQ-001/006 without using
\((\mathcal R_B)\) for those properties.
That hypothesis was used only to obtain real height
support.  Here the same support follows from the just-proved
half-plane nonvanishing and the completed functional
equation.  It reflects the divisor to the other open
half-plane; the conductor exponential has no divisor.
This includes both completed endpoints.
Thus MA-FQ-006 applies and the entire comparison datum
is zero.  Its good-prime coefficient
\(4(\log p)p^{-1/2}\) at every \(p\in\Omega\) then forces
\(\Omega\) to be empty, a contradiction.
We have proved (MA-RB.18), without a real-zero premise.

Landau's theorem at the now finite abscissa \(\gamma\)
forces a singularity of \(D_\Omega\).
Its meromorphic continuation makes that singularity
a pole.  Its positive integer order is \(k\) in
(MA-RB.19).  The finite local unit in (MA-RB.17) does
not change that order.

\emph{The symmetric-square constraint at this actual point.}
Every finite local unit and gamma unit in the proofs of
MA-S5-003 and MA-S5-006 remains holomorphic and nonzero
at any real point of \((0,1)\), not only at \(1/2\):
\(|p^{-\sigma}|<1\) excludes a finite local zero, while
the real gamma arguments there are positive.
The already proved meromorphic tensor identities therefore
give their pointwise order identities at \(\gamma\).

For clarity let \(\Pi\) and \(\nu\) be the four-dimensional
cuspidal representation and finite-order similitude
character in MA-S5-006.  If \(r_4\) is the order at \(\gamma\)
of its completed Rankin--Selberg function, then
\[
 \operatorname{End}\xi=\mathbf1\oplus V\eta\oplus W\oplus X,
 \qquad
 r_4=v_1+\operatorname{ord}_\gamma L_W+e-2k .
 \tag{MA-RB.23}
\]
The real function \(\zeta_{\mathbb Q}\) is nonzero on
\((0,1)\), by the alternating-series argument already
given in MA-RB-001.
Equation (MA-RB.17) gives
\(\operatorname{ord}_\gamma L_W=n+k\).
The exterior-square identity and tensor decomposition give
\[
 L^S(s,\Pi\times\Pi^\vee)
  =L^S(s,\Pi,\operatorname{Sym}^2\otimes\nu^{-1})
                               \zeta_{\mathbb Q}^S(s)A^S(s).
 \tag{MA-RB.24}
\]
The finite deleted factors are units at \(\gamma\).
The central character of \(\Pi\) is unitary and
\(\nu\) is finite order, as checked in MA-S5-006.

Takeda's published Theorem~7.1 states that this incomplete
twisted symmetric-square function is holomorphic everywhere
except for possible simple poles at zero and one.
In particular, it is holomorphic at the present
\(\gamma\in(1/2,1)\).
If \(q_\gamma\ge0\) is its order there, (MA-RB.23)--(MA-RB.24)
give
\[
       r_4=n+v_1+e-k=q_\gamma+n,\qquad
                         q_\gamma=v_1+e-k\ge0.
 \tag{MA-RB.25}
\]
This is an extension of the order deduction to a new real
point, not a new assertion of symmetric-square nonvanishing.
It proves the third line of (MA-RB.20).

The first two lines of (MA-RB.20) follow from \(R=A/L_W\)
and \(B=C L_W\), with signed orders retained.
Here \(n,b,v_0,v_1,e\ge0\) by standard Hecke holomorphy
and Calegari's established four- and six-dimensional
automorphy.  Also
\[
 L_{W\eta}=\zeta_F/\zeta_{\mathbb Q}
\]
is holomorphic at this particular real point; denote its
nonnegative order by \(x_1\).
Thus the actual regular-character factorization gives
\[
 \begin{aligned}
 \operatorname{ord}_\gamma\zeta_L
    &=b+4v_0+4v_1+5(n+k)+5x_1+6e\\
    &\ge b+5n+5k+4(v_1+e)
     \ge b+5n+9k .
 \end{aligned}
 \tag{MA-RB.26}
\]
No general holomorphy of the unknown \(L_W\) is being
assumed: its positive order \(n+k\) at this point is
already forced by the actual comparison pole.
Nor was \(C(\gamma)\ne0\) needed.  Its possibly positive
order \(b\) only strengthens (MA-RB.26).

\emph{The discriminant bound.}
The completed Dedekind logarithmic-derivative identity
and estimate (MA-RB.8)--(MA-RB.13) were proved from
unconditional analytic inputs; only the earlier deduction
of multiplicity from an arbitrary \(B\)-zero used
\(\mathsf A\).
They therefore apply here without that theorem.
Set \(\delta=1-\gamma\), so \(0<\delta<1/2\).
The at least \(9k\) copies of this real zero give
\[
                 \frac{9k}{t+\delta}
                     \le\frac1t+\mathcal H_L
                                  \qquad(0<t\le1).
 \tag{MA-RB.27}
\]
The optimizing value \(t=(3\sqrt{k}-1)/\mathcal H_L\)
is in the permitted interval.
Indeed \(t=1\) in (MA-RB.27) gives
\[
        9k\le(1+\delta)(1+\mathcal H_L)
                    <\tfrac32(1+\mathcal H_L).
\]
Since \(\mathcal H_L\ge1\), this implies
\(3\sqrt{k}-1<\mathcal H_L\).
At that optimizing value, solving (MA-RB.27) for
\(\delta\) yields exactly the first inequality in
(MA-RB.21).  Since \(k\ge1\), its numerator is at least
four.  This proves the remaining assertions, including
the shared-zero criterion (MA-RB.22).
\end{proof}

\paragraph{MA-RB-003 (what positivity on the rightmost line actually proves).}
In the nontrivial branch of MA-RB-002 define
\[
                    a_\gamma(t)=
                       -\operatorname{ord}_{\gamma+it}D_\Omega,
                              \qquad t\in\mathbb R .
 \tag{MA-RB.28}
\]
Then \(a_\gamma\) is an integer-valued positive-definite
function on the additive group \(\mathbb R\) equipped with
the discrete topology.  It satisfies
\[
        a_\gamma(0)=k,\qquad |a_\gamma(t)|\le k,
                         \qquad a_\gamma\in B(\mathbb R_{\rm d}).
 \tag{MA-RB.29}
\]
Its ordinary real support is locally finite.
Consequently MA-FQ-002 expresses it as a finite integer
linear combination of indicators of translated lattices
and singletons.

This is a boundary constraint.  It does not give the
real-supported explicit-formula hypothesis of MA-FQ-006.
More precisely, the actual good-prime measure at the
center \(1/2\),
\[
 \mathcal M_{\rm good}
   =4\sum_{p\in\Omega}\sum_{\substack{j\ge1\\j\ {\rm odd}}}
       (\log p)p^{-j/2}
             \bigl(\delta_{j\log p}+\delta_{-j\log p}\bigr),
 \tag{MA-RB.30}
\]
is not tempered when \(\gamma>1/2\).
The positive finite measures producing (MA-RB.29) lose
their mass at infinity in the ordinary real topology.
There is no ordinary finite-measure weak limit to which
one can transfer their exact prime-frequency support.

\begin{proof}
The logarithmic derivative in the convergence half-plane is
\[
 F(s)=-\frac{D_\Omega'(s)}{D_\Omega(s)}
       =4\sum_{p\in\Omega}\sum_{\substack{j\ge1\\j\ {\rm odd}}}
                              (\log p)p^{-js}.
 \tag{MA-RB.31}
\]
It converges absolutely for \(\operatorname{Re}s>\gamma\).
For example choose an intermediate exponent between
\(\gamma\) and \(\operatorname{Re}s\); the extra power
dominates the logarithm, and the remaining prime powers
are bounded geometric series.
For \(\varepsilon>0\), define the positive finite measure
on the ordinary real line
\[
 \nu_\varepsilon
   =4\varepsilon\sum_{p\in\Omega}
          \sum_{\substack{j\ge1\\j\ {\rm odd}}}
              (\log p)p^{-j(\gamma+\varepsilon)}
                                      \delta_{j\log p}.
 \tag{MA-RB.32}
\]
Its Fourier transform, with the unscaled minus convention,
is
\[
 \int e^{-itx}\,d\nu_\varepsilon(x)
                         =\varepsilon F(\gamma+\varepsilon+it).
 \tag{MA-RB.33}
\]
At every point \(\gamma+it\), meromorphy of \(D_\Omega\)
makes the right side tend to \(a_\gamma(t)\).
Indeed a local order \(m\) gives
\(-D_\Omega'/D_\Omega=-m/(s-\gamma-it)+O(1)\).
The same conclusion gives zero at every regular nonzero
point.

For finite \(t_1,\ldots,t_r\) and complex \(z_1,\ldots,z_r\),
\[
 \sum_{\ell,m}z_\ell\overline{z_m}
   \varepsilon F(\gamma+\varepsilon+i(t_\ell-t_m))
    =\int\left|\sum_\ell z_\ell e^{-it_\ell x}\right|^2
                                               d\nu_\varepsilon(x)\ge0 .
 \tag{MA-RB.34}
\]
Passing to the finite pointwise limit proves positive
definiteness.  Its two-by-two matrices give
\(|a_\gamma(t)|\le a_\gamma(0)=k\).
The usual GNS construction for the discrete group
represents \(a_\gamma\) as a matrix coefficient of a
unitary representation, proving membership in
\(B(\mathbb R_{\rm d})\), exactly as in MA-FQ-001.
Meromorphic local finiteness of the divisor gives local
finiteness of its support on the line
\(\operatorname{Re}s=\gamma\).
The values belong to the finite integer set
\(\{-k,\ldots,k\}\).
The already proved Cohen/coset result MA-FQ-002 now
gives the asserted finite lattice/singleton description.
No ordinary continuity of \(a_\gamma\) is asserted.

To check temperedness directly, suppose contrary to
(MA-RB.30)'s claim that its positive measure were tempered.
A fixed nonnegative compact smooth function, at least one
on \([-1,1]\), scaled to \([-T,T]\), gives
\(\mathcal M_{\rm good}([-T,T])\le C(1+T)^N\):
each relevant Schwartz seminorm of the scaled test is
bounded by a fixed power of \(1+T\), and positivity
bounds the interval mass by that pairing.
For every \(\sigma>1/2\), this polynomial mass bound implies
\[
 \int_0^\infty e^{-(\sigma-1/2)x}\,d\mathcal M_{\rm good}(x)
       =4\sum_{p\in\Omega}\sum_{j\ {\rm odd}}
                               (\log p)p^{-j\sigma}<\infty .
 \tag{MA-RB.35}
\]
This follows by summing over unit intervals, with an
exponentially decreasing majorant times a polynomial.
Since \(j\log p\ge\log2\), convergence of this positive
series also gives convergence of
\(\log D_\Omega(\sigma)
 =4\sum_{p,j\ {\rm odd}}p^{-j\sigma}/j\).
Monotone convergence of the finite Euler products then
gives convergence of its nonnegative Dirichlet series.
It follows that \(\gamma\le1/2\), contrary to MA-RB-002.
Thus \(\mathcal M_{\rm good}\) is not tempered.
Adding the known finite-total-variation bad-prime and
conductor correction cannot repair that failure.

Finally,
\[
        \nu_\varepsilon(\mathbb R)
             =\varepsilon F(\gamma+\varepsilon)\longrightarrow k>0 .
 \tag{MA-RB.36}
\]
On every fixed compact interval only finitely many
prime-power logarithms occur in (MA-RB.32), and each
coefficient tends to zero.  Thus
\(\nu_\varepsilon\) converges vaguely to zero while
its total mass tends to \(k\); it is not tight.
It has no narrow weak limit as finite measures on the ordinary
real line (that is, against bounded continuous test functions).
Equivalently its pointwise Fourier limit
\(a_\gamma\) is not ordinarily continuous at zero:
local finiteness gives a punctured neighborhood on which
it vanishes, whereas \(a_\gamma(0)=k\).
This is consistent with its positive definiteness on
\(\mathbb R_{\rm d}\), but supplies no ordinary
pure-atomic Fourier identity with the original prime
measure.  That is the precise identity used in
MA-FQ-003/006 and absent from this boundary argument.
\end{proof}

\paragraph{Tests of the remaining hypotheses.}
The simple profile \(a(t)=k\mathbf1_{\{0\}}(t)\) satisfies
every conclusion (MA-RB.29) and has locally finite support:
its finite positive-definiteness form is
\(k\sum_{\text{equal }t_\ell}|\sum z_\ell|^2\ge0\).
As an ordinary atomic measure it is \(k\delta_0\),
whose ordinary Fourier transform has a constant
continuous density.  The boundary construction has no
pure-atomic Fourier premise that would exclude this
possibility.  This is a test of the stated boundary
conditions, not an arithmetic example or a genuine
nonzero comparison \(L\)-datum.

The actual pole at \(s=\gamma>1/2\) itself violates
the real-height support premise: in the convention
\(s=1/2+iz\), it is the nonreal height
\(z=-i(\gamma-1/2)\).
The new positivity argument therefore does not extend
MA-FQ-006 by simply forgetting that premise.
Its output is (MA-RB.29), not the stronger ordinary
Fourier identity of that theorem.

The sharpened order constraints also remain formally
compatible.  At a hypothetical reflected pair
\(\gamma,1-\gamma\), assign
\[
 (b,n,v_0,v_1,e,\operatorname{ord}L_W,
                      \operatorname{ord}L_{W\eta},k)
                         =(0,0,0,1,0,1,0,1).
 \tag{MA-RB.37}
\]
Then \(q_\gamma=0\), \(r_4=0\), and
\(r_5=r_4+v_0+\operatorname{ord}L_{W\eta}=0\);
the regular order is nine.
All listed holomorphy and induced-character order
conditions are nonnegative.  Putting the same orders
at the reflected point is compatible with the
functional equations.  Unlike at the center, a
root-number-\(+1\) equation does not force the order
at one member of an off-center reflected pair to be
even.  This is only a formal order test, not a field
or a claim that these orders are arithmetically realized.

Heilbronn's simple-real-zero descent theorem still does
not apply, since the actual normal-closure multiplicity
would be at least nine.  Takeda's theorem supplies
holomorphy, not the nonvanishing of \(L_{V\eta}L_X\)
on the residual interval.  Neither a shared-zero
exclusion there nor \((\mathcal R_B)\) has been proved.
Moreover even an exact comparison would give
\(B=A C\); it would not by itself rule out real
off-center zeros of the automorphic numerator.

\paragraph{Exact sources and proof scope.}
The classical positive-Euler comparison is Calegari
(2013), Theorem~4.7, article-PDF p.~12, Lemmas~4.8--4.10,
article-PDF pp.~13--15.  These are the PDF page numbers
used in MT-S508, not the journal pagination.
The finite-factor normalization and
Landau application were checked in MT-S508 and MA-S5-001.
The actual gamma equality and genuine comparison
\(L\)-datum inputs are MA-FQ-001/006.
The finite representation identities and bad-local
unit assertions are the already checked MA-S5-003/006;
their certificates are not rerun here.

The new pointwise symmetric-square application uses
Takeda (2015), Theorem~7.1 and its entire terminal proof,
primary preprint \texttt{arXiv:1506.04791v1},
printed pp.~49--50.
Those two pages were freshly read to verify the
full holomorphy range used in (MA-RB.25).
Its earlier Eisenstein-series and integral-representation
theorems remain the same declared classical inputs
as in MA-S5-006; no new audit of all their proofs is
claimed.
The required Hadamard constant, gamma signs and
Dedekind estimate are proved in MA-RB-001,
(MA-RB.8)--(MA-RB.13), independently of its separate
use of \(\mathsf A\) to convert an arbitrary \(B\)-zero.
The boundary Fourier-series and non-temperedness
arguments are proved in full above.

The coefficient nine in (MA-RB.26) combines the
regular Artin-motive multiplicity five with the
symmetric-power identity
\((\operatorname{Sym}^2\xi)\otimes\nu^{-1}
=V\eta\oplus X\) already placed in the representation
and motive framework of MA-S5.
The exact inherited locators remain Deligne (1979),
Definition~6.1 and \S\S6.2--6.3, and SGA~1,
Expos\'e~V, \S7.
This interpretation accounts for the algebraic factors;
it does not supply the missing real nonvanishing.

Takeda, S. (2015). On a certain metaplectic Eisenstein series
and the twisted symmetric square \(L\)-function.
\emph{Mathematische Zeitschrift, 281}, 103--157.
\url{https://doi.org/10.1007/s00209-015-1476-x}.
